% ============================================
% Johns Hopkins Letter Template
% Uses the built-in 'letter' class
% ============================================

\documentclass[11pt,letterpaper]{letter}

\usepackage{graphicx}
\usepackage{eso-pic}
\usepackage{geometry}
\usepackage{microtype}
\usepackage[utf8]{inputenc}
\usepackage[hidelinks]{hyperref}

% ----- Page Setup -----
\geometry{left=25mm,right=25mm,top=25mm,bottom=25mm}

% ----- Logo Placement (foreground, first page only) -----
\newcommand{\PutJHULogoFG}[1]{%
  \AddToShipoutPictureFG*{%
    \AtPageUpperLeft{%
      \hspace{10mm}%
      \raisebox{-38mm}{%
        \includegraphics[width=6cm]{#1}%
      }%
    }%
  }%
}

% ----- Sender Information -----
\signature{Decory Edwards}
\address{Department of Economics \\ Johns Hopkins University \\ Baltimore, MD 21218 \\ \texttt{dedwar65@jh.edu}}

% ----- Recipient Info -----
\begin{document}
\PutJHULogoFG{Images/jhulogo.png}

\begin{letter}{Hiring Committee \\
CFM \\
New York, NY 10001}

\opening{Dear Hiring Committee,}

I am writing to express my interest in the Quantitative Researcher position at CFM. I am a Ph.D.\ candidate in the Department of Economics at Johns Hopkins University (advisors: Christopher D.\ Carroll, Nicholas W.\ Papageorge, and Stelios Fourakis), and I expect to receive my degree in May 2026. My primary specializations are macroeconomic modeling, wealth and income dynamics, and computational economics.

My research agenda focuses on the ability of macroeconomic models to generate empirically realistic distributions of wealth and income over the life cycle. A key driver of that work is modeling heterogeneity in the rate of return on assets, and understanding how return dispersion contributes to portfolio accumulation across households.

In my job market paper, I calibrate a life cycle heterogeneous agent model to match wealth moments from the Survey of Consumer Finances by allowing households to differ in their portfolio returns. The model helps explain observed wealth concentration and return dispersion. In a related research project using the Health and Retirement Study, I examine how trust influences both income growth and asset returns. I find a hump shaped relationship for returns and characterize the distribution of the persistent component of returns to net worth. These experiences have equipped me to frame research strategies and design modeling approaches, execute econometric and simulation analysis with large micro data sets, and translate results into clear and rigorous written deliverables for both academic and practitioner audiences.

I am particularly interested in CFM because of its global and systematic approach to asset management, where researchers work within a unified research platform rather than narrow product specific teams. A central component of CFM’s work involves identifying predictors of returns, testing their stability across asset classes and market environments, and integrating them into portfolios subject to explicit risk constraints. My training aligns directly with this process. In my research, I routinely evaluate whether estimated return differences reflect persistent structure or noise, assess robustness across samples, and interpret model outputs through an economic lens. The emphasis at CFM on careful empirical validation, large scale data analysis, and close collaboration between researchers and portfolio managers closely mirrors how I have been trained to conduct quantitative research.

I am enthusiastic about the opportunity to live and work in New York and to engage fully in the in person research environment that CFM expects. I am comfortable working on site according to the firm’s schedule and value the collaboration that comes from regular in person interaction. I am also drawn to the prospect of a three year role, which would allow me to deepen my applied research skills while continuing to develop my broader research agenda in an environment where rigorous empirical analysis directly informs real world investment decisions.

I believe I would be a strong fit at CFM because:
\begin{enumerate}
    \item My experience with simulation based modeling and empirical validation aligns with research on return forecasting and signal evaluation.
    \item My work on heterogeneity in returns prepares me to analyze variation in asset behavior across market conditions.
    \item I bring strong communication skills and a collaborative approach that supports effective interaction between research and portfolio teams.
\end{enumerate}

I would welcome the opportunity to contribute my quantitative training, modeling experience, and interest in applied research to CFM. Thank you for considering my application.

\closing{Sincerely,}

\end{letter}
\end{document}
